% Устойчивость к замене языковой модели: основная и локальная малая модель, проценты.
% Значения: eval/reports/2026-11-11-comparison.md (deepseek-chat),
% eval/reports/2026-11-17-hybrid-qwen4b.md и 2026-11-17-llm_only-qwen4b.md (qwen3.5-4b).
\begin{tikzpicture}[
  font=\footnotesize,
  x=1cm, y=0.052cm,
  bar/.style={draw, semithick},
  val/.style={font=\scriptsize, above, inner sep=1pt},
]

\pgfkeys{/pgf/number format/.cd, use comma, fixed, fixed zerofill, precision=1}

% --- оси и сетка ---
\foreach \v in {0, 20, 40, 60, 80, 100} {
  \draw[very thin, gray!60] (0, \v) -- (13.4, \v);
  \node[left] at (0, \v) {\v};
}
\draw[semithick] (0, 0) -- (13.4, 0);
\draw[semithick] (0, 0) -- (0, 104);
\node[rotate=90] at (-0.95, 50) {Доля, \%};

% --- группы столбцов: только модель, гибрид ---
\foreach \x/\a/\c/\name in {
  0.75/62.3/91.8/{Совпало\\с разметкой,\\deepseek-chat},
  4.05/36.1/85.2/{Совпало\\с разметкой,\\qwen3.5-4b},
  7.35/26.2/0.0/{Недопустимые\\действия,\\deepseek-chat},
  10.65/47.5/0.0/{Недопустимые\\действия,\\qwen3.5-4b}} {
  \draw[bar, fill=white] (\x, 0) rectangle ++(0.9, \a);
  \node[val] at (\x + 0.45, \a) {\pgfmathprintnumber{\a}};
  \draw[bar, fill=black!75] (\x + 1.0, 0) rectangle ++(0.9, \c);
  \node[val] at (\x + 1.45, \c) {\pgfmathprintnumber{\c}};
  \node[align=center, below] at (\x + 0.95, -2) {\name};
}

% --- условные обозначения ---
\draw[bar, fill=white] (3.0, -46) rectangle ++(0.5, 7);
\node[right] at (3.55, -42.5) {только языковая модель};
\draw[bar, fill=black!75] (8.4, -46) rectangle ++(0.5, 7);
\node[right] at (8.95, -42.5) {гибрид};

\end{tikzpicture}
